\chapter{Unit 3: Orthographic Projections of Objects}

\section{A. What You'll Learn}
In this unit, you transition from points and lines to complete 3D mechanical components and objects. You will master:
\begin{itemize}
    \item Multi-view orthographic projection theory and the **Glass Box Concept**.
    \item ISO First Angle vs. ANSI Third Angle view arrangements and symbol differentiation.
    \item The **Miter-Line Method ($45^\circ$)** for depth transfer between Top View and Side Views.
    \item Application of visible object lines, hidden lines, and center lines on engineering views.
    \item AutoCAD Layer Management (`LINETYPE`, `LTSCALE`).
    \item AutoCAD 2D Modification commands (`MOVE`, `COPY`, `ROTATE`, `TRIM`, `MIRROR`, `SCALE`, `FILLET`, `CHAMFER`, `ARRAY`).
    \item Complete step-by-step AutoCAD 2D drafting workflow for multi-view orthographic drawings.
\end{itemize}

---

\section{B. Multi-View Orthographic Projection Principles}

\subsection{The "Glass Box" Concept}
Imagine a 3D object suspended inside a transparent glass box.
\begin{enumerate}
    \item The observer looks perpendicular ($90^\circ$) to each face of the box.
    \item Points and features of the object are projected onto the transparent glass walls.
    \item The glass box is unfolded flat into a single 2D plane surface.
    \item The resulting 2D drawings form the set of orthographic views: **Front View (FV)**, **Top View (TV)**, **Left Side View (LSV)**, and **Right Side View (RSV)**.
\end{enumerate}

\subsection{View Layout Systems: First Angle vs. Third Angle}

\begin{importantrule}[First Angle View Placement (ISO Standard)]
The object is between Observer and Plane of Projection.
\begin{itemize}
    \item **Front View (FV):** Placed above the $XY$ reference line.
    \item **Top View (Plan):** Placed directly **below** the Front View.
    \item **Left Side View (LSV):** Drawn to the **right** of the Front View.
    \item **Right Side View (RSV):** Drawn to the **left** of the Front View.
\end{itemize}
\end{importantrule}

\begin{importantrule}[Third Angle View Placement (ANSI Standard)]
The Plane of Projection is between Observer and Object.
\begin{itemize}
    \item **Front View (FV):** Placed below the $XY$ reference line.
    \item **Top View (Plan):** Placed directly **above** the Front View.
    \item **Right Side View (RSV):** Drawn to the **right** of the Front View.
    \item **Left Side View (LSV):** Drawn to the **left** of the Front View.
\end{itemize}
\end{importantrule}

---

\section{C. The Miter-Line Method for View Projection}

When constructing 2D orthographic projections manually or in AutoCAD, dimensions are correlated across views:
\begin{itemize}
    \item \textbf{Front View \& Top View:} Share identical \textbf{Length (Width)} dimensions horizontally.
    \item \textbf{Front View \& Side View:} Share identical \textbf{Height} dimensions vertically.
    \item \textbf{Top View \& Side View:} Share identical \textbf{Depth} dimensions.
\end{itemize}

\begin{keyequation}[Miter Line ($45^\circ$) Depth Transfer]
To project depth dimensions accurately between Top View and Side View:
\begin{enumerate}
    \item Draw a diagonal construction line at $45^\circ$ starting from the intersection of the main reference axes in the projection quadrant.
    \item Project horizontal lines from Top View features to meet the $45^\circ$ Miter Line.
    \item At each intersection point on the Miter Line, project vertical lines down (or up) into the Side View area.
    \item Project horizontal lines from Front View across into the Side View. The grid intersection forms the Side View geometry.
\end{enumerate}
\end{keyequation}

---

\section{D. AutoCAD Layer Properties \& Line Management}

Technical drawings rely on standardized line properties managed via the **Layer Properties Manager**.

\begin{table}[htbp]
\centering
\small
\begin{tabularx}{\textwidth}{l c c X}
\toprule
\textbf{Layer Name} & \textbf{Color} & \textbf{Line Type} & \textbf{Line Weight \& Purpose} \\
\midrule
\textbf{Object Lines} & White / Black & Continuous & $0.5\text{--}0.7\text{ mm}$ (Thick visible outlines) \\
\textbf{Hidden Lines} & Yellow & \texttt{DASHED} & $0.3\text{ mm}$ (Medium interior hidden details) \\
\textbf{Center Lines} & Red & \texttt{CENTER} & $0.3\text{ mm}$ (Thin long-short-long axes) \\
\textbf{Projectors} & Grey / Cyan & Continuous & $0.15\text{ mm}$ (Thin layout guide lines) \\
\textbf{Dimensions} & Green & Continuous & $0.25\text{ mm}$ (Dimension text and arrows) \\
\bottomrule
\end{tabularx}
\caption{Standard AutoCAD Layer Property Configuration.}
\label{tab:autocad-layers}
\end{table}

\begin{autocadcmd}[LTSCALE]
Globally scales the length of dashes and gaps in non-continuous line types (e.g., \texttt{DASHED}, \texttt{CENTER}).
\begin{itemize}
    \item \textbf{Command:} \texttt{LTSCALE}
    \item \textbf{Usage:} If dashed center lines appear as continuous solid lines, type \texttt{LTSCALE} $\to$ Enter scale factor (e.g., \texttt{0.5} or \texttt{10}).
\end{itemize}
\end{autocadcmd}

---

\section{E. AutoCAD 2D Modification Commands}

\begin{autocadcmd}[MOVE \& COPY]
Displaces or duplicates objects.
\begin{itemize}
    \item \textbf{Move Command:} \texttt{MOVE} or \texttt{M}
    \item \textbf{Copy Command:} \texttt{COPY} or \texttt{CO} / \texttt{CP}
    \item \textbf{Workflow:} Select objects $\to$ Press \texttt{Enter} $\to$ Specify Base Point $\to$ Specify Second Point.
\end{itemize}
\end{autocadcmd}

\begin{autocadcmd}[ROTATE]
Revolves selected objects around a base point.
\begin{itemize}
    \item \textbf{Command:} \texttt{ROTATE} or \texttt{RO}
    \item \textbf{Direction:} Positive angle values rotate **Counter-Clockwise**; negative values rotate Clockwise.
\end{itemize}
\end{autocadcmd}

\begin{autocadcmd}[TRIM]
Cuts objects to meet the edges of other boundary objects.
\begin{itemize}
    \item \textbf{Command:} \texttt{TRIM} or \texttt{TR}
    \item \textbf{Workflow:} Select cutting edges (or press \texttt{Enter} for select-all) $\to$ Click object segments to remove.
\end{itemize}
\end{autocadcmd}

\begin{autocadcmd}[MIRROR]
Creates a reverse symmetrical copy of selected objects across a defined axis.
\begin{itemize}
    \item \textbf{Command:} \texttt{MIRROR} or \texttt{MI}
    \item \textbf{Workflow:} Select objects $\to$ Specify 1st and 2nd points of mirror line $\to$ Choose whether to delete source objects (\texttt{Yes/No}).
\end{itemize}
\end{autocadcmd}

\begin{autocadcmd}[FILLET \& CHAMFER]
Rounds or bevels corner edges of objects.
\begin{itemize}
    \item \textbf{Fillet Command:} \texttt{FILLET} or \texttt{F} (Set radius first: Type \texttt{F} $\to$ \texttt{R} $\to$ Enter radius value $\to$ Select two lines).
    \item \textbf{Chamfer Command:} \texttt{CHAMFER} or \texttt{CHA} (Set distance first: Type \texttt{CHA} $\to$ \texttt{D} $\to$ Enter Dist 1 \& Dist 2 $\to$ Select two lines).
\end{itemize}
\end{autocadcmd}

\begin{autocadcmd}[ARRAY]
Creates multiple copies of objects in a rectangular, polar, or path pattern.
\begin{itemize}
    \item \textbf{Command:} \texttt{ARRAY} or \texttt{AR}
    \item \textbf{Rectangular Array:} Arranges copies into Rows and Columns grid.
    \item \textbf{Polar Array:} Copies objects in a circular pattern around a center axis point.
\end{itemize}
\end{autocadcmd}

---

\section{F. Hands-on AutoCAD Workflow: 2D Orthographic Views}

Follow this exact multi-step procedure when converting a 3D isometric view into 2D orthographic drawings:

\begin{enumerate}
    \item \textbf{Setup Phase:}
    \begin{itemize}
        \item Set units: Type \texttt{UN} $\to$ Decimal, $0.00$ precision, Millimeters.
        \item Set limits: Type \texttt{LIMITS} $\to$ \texttt{0,0} to \texttt{210,297} $\to$ Type \texttt{Z} $\to$ \texttt{A}.
        \item Create Layers: Object Lines (Thick White), Hidden Lines (Yellow Dashed), Center Lines (Red Center), Projectors (Grey).
    \end{itemize}
    \item **Construction Grid:**
    \begin{itemize}
        \item Draw horizontal and vertical reference lines using \texttt{XLINE} or \texttt{LINE}.
        \item Enable \texttt{ORTHO (F8)} for straight projection lines.
    \end{itemize}
    \item **Front View Construction:**
    \begin{itemize}
        \item Switch to 'Object Layer'.
        \item Draw Front View outline using \texttt{LINE}, \texttt{CIRCLE}, \texttt{FILLET}.
    \end{itemize}
    \item **Top View Construction:**
    \begin{itemize}
        \item Switch to 'Projector Layer'. Use \texttt{XLINE} to project vertical lines down from Front View corners.
        \item Switch to 'Object Layer' and draw Top View width and depth outlines.
        \item Add hidden lines for internal holes using 'Hidden Layer'.
    \end{itemize}
    \item **Side View Construction (Miter Line Method):**
    \begin{itemize}
        \item Draw a $45^\circ$ line from reference axis intersection.
        \item Project horizontal lines from Top View rightwards to meet $45^\circ$ line.
        \item Project vertical lines down from $45^\circ$ line into Side View quadrant.
        \item Project horizontal lines from Front View into Side View quadrant.
        \item Trace Side View outlines at grid intersections.
    \end{itemize}
    \item \textbf{Final Details \& Dimensioning:}
    \begin{itemize}
        \item Mark hole centers using 'Center Layer'.
        \item Add linear dimensions (\texttt{DLI}) and radial dimensions (\texttt{DRA}).
    \end{itemize}
\end{enumerate}

---

\section{G. Chapter 3 Practice Problems \& MCQs}

\subsection{Practice Questions}

\begin{vivaqa}[Question 1 (Basic)]
Explain the "Glass Box Concept" in orthographic projections.\\[4pt]
\textbf{Answer:} The 3D object is imagined suspended inside a transparent glass box. Views are projected perpendicularly onto the faces of the box, which is then unfolded flat into a 2D plane.
\end{vivaqa}

\begin{vivaqa}[Question 2 (Moderate)]
In First Angle Projection, where are the Top View and Left Side View placed relative to the Front View?\\[4pt]
\textbf{Answer:} Top View is placed **below** the Front View. Left Side View is drawn to the **right** of the Front View.
\end{vivaqa}

\begin{vivaqa}[Question 3 (Exam Level)]
Describe the steps to draw a Fillet of radius $15\text{ mm}$ on a sharp corner in AutoCAD.\\[4pt]
\textbf{Answer:} Type \texttt{FILLET} (or \texttt{F}) $\to$ Press \texttt{Enter} $\to$ Type \texttt{R} (Radius) $\to$ Press \texttt{Enter} $\to$ Type \texttt{15} $\to$ Press \texttt{Enter} $\to$ Click first line $\to$ Click second line.
\end{vivaqa}

\subsection{Multiple-Choice Questions (MCQs)}

1. In First Angle Projection, the Right Side View is drawn to the:  
   A. Right of Front View  
   B. Left of Front View  
   C. Top of Front View  
   D. Bottom of Front View  
   \textbf{Answer: B} — Right Side View is drawn to the left of Front View in First Angle Projection.

2. Which AutoCAD command is used to round sharp corners with a specific radius?  
   A. \texttt{CHAMFER}  
   B. \texttt{FILLET}  
   C. \texttt{ARC}  
   D. \texttt{CIRCLE}  
   \textbf{Answer: B} — \texttt{FILLET} rounds sharp corners.

3. What angle is used for the Miter Line to transfer depth between Top View and Side View?  
   A. $30^\circ$  
   B. $45^\circ$  
   C. $60^\circ$  
   D. $90^\circ$  
   \textbf{Answer: B} — The Miter Line is drawn at $45^\circ$.

4. Which line type is used to represent hidden interior edges in orthographic views?  
   A. Continuous Thick Line  
   B. Long dash -- short dash Chain Line  
   C. Dashed Medium Line  
   D. Continuous Thin Line  
   \textbf{Answer: C} — Dashed lines represent hidden interior edges.

5. What does the AutoCAD command \texttt{LTSCALE} control?  
   A. Scale of drawing limits  
   B. Global dash and gap spacing of non-continuous lines  
   C. Text height in dimension style  
   D. Zoom magnification level  
   \textbf{Answer: B} — \texttt{LTSCALE} controls global non-continuous line dash scale.

---

\section{H. Unit 3 Quick Revision Checklist}
\begin{revisionchecklist}
\begin{itemize}
    \item $\text{First Angle: } \text{FV above } XY, \text{ TV below FV}, \text{ LSV to right of FV}, \text{ RSV to left of FV}$.
    \item $\text{Third Angle: } \text{FV below } XY, \text{ TV above FV}, \text{ RSV to right of FV}, \text{ LSV to left of FV}$.
    \item $\text{Depth Transfer: } \text{Use } 45^\circ \text{ Miter Line between TV and Side View}$.
    \item $\text{Visible Outlines } \to 0.5\text{--}0.7\text{ mm Thick}$; $\text{Hidden Lines } \to 0.3\text{ mm Dashed}$; $\text{Center Lines } \to \text{Chain Line}$.
    \item $\text{AutoCAD Fillet: } \texttt{F} \to \texttt{R} \to \text{Radius Value} \to \text{Select Lines}$.
    \item $\text{AutoCAD Chamfer: } \texttt{CHA} \to \texttt{D} \to \text{Dist 1, Dist 2} \to \text{Select Lines}$.
    \item $\text{AutoCAD Rotate: } \text{Positive angle rotates Counter-Clockwise}$.
    \item $\text{AutoCAD Array: } \text{Rectangular (Rows/Cols) vs. Polar (Circular)}$.
\end{itemize}
\end{revisionchecklist}
